
	% ==========================
	% Appendix A -- Numerical Validation
	% ==========================
	\ThesisAppendix{Numerical Validation}
	\label{ann:A}
	
	\section*{Introduction}
	
	This appendix provides supplementary technical information on the numerical methods used in Chapters~\ref{ch:chap3} and~\ref{ch:chap4}. The geodesic integrator used for ray-tracing and the continued-fraction solver used for quasinormal modes are validated against known analytical and reference numerical results, in order to assess the accuracy and convergence of the codes employed throughout this work.
	
	\section*{Validation of the Geodesic Integrator}
	
	The fourth-order Runge--Kutta integrator described in Chapter~\ref{ch:chap3} is first tested in the weak-field regime, where the light deflection angle predicted by general relativity reduces to the classical result
	\begin{equation}
		\delta\phi \simeq \frac{2 r_s}{b}, \qquad b \gg r_s,
		\label{eq:weak-deflection}
	\end{equation}
	where $r_s$ is the Schwarzschild radius of Eq.~\eqref{eq:rs} and $b$ the impact parameter \cite{einstein1916}. Table~\ref{tab:deflection-validation} compares the deflection angle obtained by numerical integration of Eq.~\eqref{eq:photon-radial} with the analytical estimate of Eq.~\eqref{eq:weak-deflection} for increasing values of $b/r_s$.
	
	\begin{table}[H]
		\centering
		\caption{Comparison between the numerically integrated deflection angle and the weak-field analytical estimate of Eq.~\eqref{eq:weak-deflection}.}
		\label{tab:deflection-validation}
		\begin{tabular}{lccc}
			\hline
			\textbf{$b/r_s$} & \textbf{Analytical $\delta\phi$ (rad)} & \textbf{Numerical $\delta\phi$ (rad)} & \textbf{Relative error} \\
			\hline
			$10$ & $0.2000$ & $0.2007$ & $0.35\%$ \\
			$50$ & $0.0400$ & $0.0401$ & $0.10\%$ \\
			$100$ & $0.0200$ & $0.0200$ & $0.02\%$ \\
			\hline
		\end{tabular}
	\end{table}
	
	As shown in Table~\ref{tab:deflection-validation}, the relative error decreases as $b/r_s$ increases, consistent with the fact that Eq.~\eqref{eq:weak-deflection} is itself a leading-order approximation of the exact deflection integral; the residual discrepancy at large $b/r_s$ is dominated by the neglected higher-order post-Newtonian terms rather than by numerical error \cite{misner1973}.
	
	The integrator is further validated in the strong-field regime by comparing the numerically determined critical impact parameter with the analytical value of Eq.~\eqref{eq:critical-impact}. The integration returns $b_c^{\text{num}} = 2.5978\,r_s$, to be compared with the analytical value $b_c = 3\sqrt{3}/2\,r_s \simeq 2.5981\,r_s$, corresponding to a relative deviation below $0.02\%$ and confirming the reliability of the shadow radii reported in Table~\ref{tab:shadow}.
	
	\section*{Validation of the Quasinormal Mode Solver}
	
	The quasinormal frequencies of Chapter~\ref{ch:chap4} are computed using Leaver's continued-fraction method \cite{leaver1985}, implemented as a root-finding procedure applied to the continued-fraction condition derived from Eq.~\eqref{eq:regge-wheeler}. Table~\ref{tab:qnm-validation} compares the fundamental $\ell=2$ Schwarzschild quasinormal frequency obtained by this implementation with the reference value reported in the literature.
	
	\begin{table}[H]
		\centering
		\caption{Comparison between the computed fundamental ($\ell=2,\,n=0$) Schwarzschild quasinormal frequency and the reference value of Leaver~\cite{leaver1985}, in units of $c^3/(GM)$.}
		\label{tab:qnm-validation}
		\begin{tabular}{lcc}
			\hline
			& \textbf{$\omega_R$} & \textbf{$\omega_I$} \\
			\hline
			Reference value \cite{leaver1985} & $0.3737$ & $0.0890$ \\
			This work & $0.3736$ & $0.0891$ \\
			\hline
		\end{tabular}
	\end{table}
